Robotic visual inspection of complex industrial components requires precise alignment between the end-effector and local surface geometry so that imaging optics remain normal to and focused on the region of interest. In practice, this alignment must be maintained in the presence of depth-sensing noise, surface irregularities, and—critically—the continuous involvement of a human inspector. Although automated systems are increasingly capable of part recognition, search, and detection, human operators remain essential for the classification and action stages of anomaly response. Existing surface-following pipelines, however, are typically organized around offline coverage planning over reconstructed geometry, or around discrete perception–correction loops driven by proximity or infrared sensing. These designs are poorly suited to shared-autonomy operation, where teleoperation inputs and perceptual updates arrive concurrently and must be reconciled in real time.
This paper presents an admittance-based surface alignment framework that unifies perception-driven orientation regulation and operator input within a single continuous control law. The end-effector is modeled as a virtual sphere of mass m and radius R moving through a viscous medium, producing a physically interpretable mass–damper system whose translational and rotational responses are deliberately synchronized through a modified rotational drag coefficient. Surface normals are estimated continuously from an eye-in-hand Intel RealSense D405 depth stream, using RANSAC-based plane fitting to reject outliers followed by principal component analysis on the resulting inlier set. The orientation error between the estimated surface normal and the camera optical axis is passed to a PD controller, whose output is interpreted as a virtual torque applied at a pivot offset from the sphere center. Newton–Euler analysis converts this torque into an equivalent force and torque at the sphere, which serve as the inputs to the admittance dynamics; integration of the resulting accelerations yields Cartesian velocity commands transmitted to the robot through the ROS 2 Servo node.
Controller gains are selected by matching the closed-loop characteristic polynomial to a critically damped second-order form, and a torque-saturation constraint derived from the actuator velocity limit and viscous drag is used to bound the admissible natural frequency for a given maximum expected orientation error. Because the framework operates as an outer-loop compliance layer on top of the manipulator's existing joint servo loops, it is compatible with standard position- or velocity-controlled industrial platforms and does not require torque-level access.
The framework was validated on a UR5e manipulator under two initial conditions: a 15° error case operating entirely within the unsaturated regime, and a 40° error case that drives the controller into torque saturation during the initial transient. In both cases, the orientation error, angular and linear velocities, and torque commands measured on hardware closely track the predictions of a synthetic MATLAB simulation of the same control pipeline, with residual deviations attributable to depth-sensing noise, surface-normal estimation uncertainty, and the internal dynamics of the manufacturer's joint controllers. The framework attains a final mean orientation error of 0.007 rad, matching the performance reported for prior position-based iterative adaptive methods that rely on infrared sensing, while additionally providing continuous closed-loop control and a force-based formulation naturally suited to human-in-the-loop operation. A baseline comparison against ten manual teleoperation trials confirms that the proposed controller delivers substantially more repeatable convergence than unassisted operator alignment.
The principal contributions of this work are: (i) a perception-driven, closed-loop orientation regulation strategy using surface normal estimation and a virtual mass–damper outer-loop compliance layer; (ii) an admittance-based formulation that maps orientation error, teleoperation inputs, and human corrections to compliant motion while respecting actuator constraints; (iii) a modular architecture compatible with industrial robots under position or velocity control; and (iv) experimental validation demonstrating stable convergence and robustness to perception noise.